\chapter{Diagramas de Bloques de Confiabilidad (RBDs)}
\label{chapter:rbd}
En este capítulo se introduce el modelo matemático sobre el cual yace este trabajo. En \ref{sec:rbd_definicion} se establece la definición formal de \gls{RBD} (\textit{Reliability Block Diagram}). En \ref{sec:rbd_semantica} se describen definiciones de carácter semántico, como \textit{path} y \textit{cut set}. Luego, en \ref{sec:rbd_composiciones} se determinan los operadores del modelo: composición en serie y en paralelo.
Por último, se enuncia y demuestra una proposición que asegura que el modelo es composicional. Adicionalmente,
a lo largo del capítulo, se justifican algunas decisiones de diseño y restricciones importantes del modelo.
Posteriormente, en el siguiente capítulo, se definirá, en términos de RBDs, los RBD coloreados, que permiten extender la expresividad del modelo para que, además de composicional, sea cerrado bajo operadores y modular.

Previo al desarrollo de la definición formal de RBD, veamos los antecedentes teóricos necesarios.
\section{Preliminares}
\label{sec:rbd_prelim}




A \textit{directed graph} $G$ is a tuple $(V,E)$ where $V$ is a set of vertices and $E \subseteq V \times V$ is a set of ordered pairs. Given a graph $\graph$, a \textit{walk} is a sequence of vertices $\walk$ such that $\forall i: 1 \le i \le n-1 : (v_i, v_{i+1}) \in E$. If all vertices are distinct, a walk is called a \textit{path}. A \textit{cycle} is a walk where the only repeated vertices are $v_1=v_k$. A \textit{DAG (Directed Acyclic Graph)} is a directed graph that has no cycles.

Given a graph $\graph$ and a subset of nodes $X \subseteq V$, the \textit{induced graph} $\indgraph{X}$ is defined as: 
%
\[\indgraph{X} = (X, E \cap (X \times X)).\]



\section{RBD: Definición formal}
\label{sec:rbd_definicion}
\completar{Definición de RBD como grafo dirigido acíclico (DAG) con nodos terminales \start y \nd.}

\section{Semántica}
\label{sec:rbd_semantica}
\completar{Semántica de conectividad, estados de los componentes y función de estructura / probabilidad de supervivencia.}

\section{Composiciones}
\label{sec:rbd_composiciones}
\completar{Operaciones de composición (serie, paralelo, puente / grafos arbitrarios).}